\documentclass[11pt,a4paper]{article}

\usepackage[utf8]{inputenc}
\usepackage[T1]{fontenc}
\usepackage[margin=2.3cm]{geometry}
\usepackage{hyperref}
\usepackage{microtype}
\usepackage{parskip}
\usepackage{lmodern}
\usepackage{enumitem}
\usepackage{titlesec}

\hypersetup{colorlinks=true, urlcolor=blue!60!black}
\titleformat{\section}{\large\bfseries}{}{0pt}{}
\titlespacing{\section}{0pt}{10pt}{4pt}
\pagestyle{empty}

\begin{document}

\begin{center}
{\large\textbf{Programming, Data-Analysis, and Machine-Learning Projects}}\\[2pt]
Kundai Farai Sachikonye \quad--\quad
\href{https://github.com/fullscreen-triangle}{github.com/fullscreen-triangle}
\end{center}

\vspace{4pt}

An overview of representative projects, grouped by the three data modalities of
xMDT-HPB --- clinical/language, imaging, and laboratory data --- followed by the ML
engineering and reproducibility work underneath them. All projects are open-source;
most are deployed with a live demo. Where a number is given it is validated against
a named reference (IEDB, NIST, 1000~Genomes/gnomAD/UK~Biobank).

\section{Explainable / auditable diagnostic AI}

\textbf{syndrome} --- \href{https://syndrome-seven.vercel.app/tools}{syndrome-seven.vercel.app}\\
Computational-oncology toolkit: neoantigen identification, MHC binding prediction
(AUC~0.81, validated on 2{,}847 IEDB peptides), and immune-escape scoring. The
method uses three-channel discrete Fourier transforms on amino-acid physicochemical
properties with \emph{zero} free parameters trained on patient data --- every
prediction is reproducible from the algorithm alone, independent of any training
cohort. Developed with KRAS-mutant pancreatic ductal adenocarcinoma (PDAC) immune
evasion as a worked clinical case. Runs in-browser (WebGPU), 50\,MB, no cluster
required. \emph{Relevance:} a concrete instance of \emph{structural} explainability
(auditable by construction) as distinct from post-hoc explanation (SHAP/LIME).

\textbf{crown-prince} --- \href{https://github.com/fullscreen-triangle/crown-prince}{github.com/fullscreen-triangle/crown-prince}\\
Pharmacological AI: PK/PD/ADME trajectories and binding affinities derived from
molecular geometry, validated 39/39 against NIST with zero empirical parameters.
Interpretable-by-construction: predictions are geometric operations, not statistical
associations from potentially biased training data.

\section{Imaging}

\textbf{hieronymus} --- \href{https://hieronymus-omega.vercel.app/}{hieronymus-omega.vercel.app}\\
GPU-accelerated medical image analysis: segmentation, point-spread-function
deconvolution, registration, and morphometry, with metric-preserving reconstruction
and uncertainty quantification. Implemented in WebGL/GLSL for real-time pipelines.
\emph{Relevance:} direct experience with the image-analysis primitives underlying
radiological AI --- segmentation, registration, quantitative morphometry.

\textbf{BRUT (rPPG)} --- \href{https://remote-photoplethysmography.vercel.app/}{remote-photoplethysmography.vercel.app}\\
Camera-based physiological inference: recovers melanin, haemoglobin, and vasodilation
from ordinary RGB video by inverting a layered optical skin model. \emph{Relevance:}
multimodal signal fusion --- extracting several physically distinct physiological
channels from one imaging stream, with an explicit forward model rather than a
black box.

\section{Laboratory \& multi-omics data}

\textbf{lavoisier} --- \href{https://lavoisier.vercel.app/sandbox}{lavoisier.vercel.app/sandbox}\\
High-throughput mass-spectrometry pipeline: platform-invariant feature extraction,
98.9\% accuracy against NIST reference libraries, distributed processing (Ray),
memory-mapped I/O for large mzML datasets. Deterministic and benchmarked ---
reproducible by a third party from a stated seed. \emph{Relevance:} the reproducible,
reference-validated laboratory-data pipeline discipline the project asks for.

\textbf{gospel} --- \href{https://github.com/fullscreen-triangle/gospel}{github.com/fullscreen-triangle/gospel}\\
Multi-site multi-omics integration: host genotype, transcriptomics, and metagenomics
as a joint Bayesian graph; validated against 1000~Genomes, gnomAD, and UK~Biobank.
\emph{Relevance:} hands-on experience with what data actually flows across multiple
clinical sites, where pseudonymisation fails, and the constraints multi-site
analysis places on data governance.

\section{ML engineering, agentic systems, and reproducibility}

\textbf{kwasa-kwasa / Turbulance} --- \href{https://github.com/fullscreen-triangle/kwasa-kwasa}{github.com/fullscreen-triangle/kwasa-kwasa}\\
A declarative orchestration language for multi-agent, multi-substrate delegation
(LLM oracles, tools, Python) with confidence/argumentation semantics and grounding
of generated values. Rust core + TypeScript, with a VSCode language server.

\textbf{four-sided-triangle} --- \href{https://github.com/fullscreen-triangle/four-sided-triangle}{github.com/fullscreen-triangle/four-sided-triangle}\\
An 8-stage metacognitive multi-agent RAG / reasoning pipeline with resource-aware
orchestration; Rust core, FastAPI + Ray + Docker.

\textbf{purpose} --- \href{https://github.com/fullscreen-triangle/purpose}{github.com/fullscreen-triangle/purpose}\\
A Rust CLI for agent-queryable context retrieval that reduces LLM token cost;
trait-based architecture in which hand-coded and LoRA-trained resolvers swap
transparently.

\vspace{6pt}
\textbf{Foundation-model adaptation \& training methodology.} LoRA fine-tuning and
knowledge distillation; research on curriculum design (inverse curricula) for
sample-efficient training. PyTorch throughout, at HPC/SLURM scale.

\vspace{6pt}
\textbf{Reproducibility discipline (across all projects).} Public source; validation
suites that are deterministic and seeded, so results reproduce exactly; every
quantitative claim checked against an explicit external reference rather than
asserted. Version control, documentation, and containerised deployment as standard.

\end{document}
